\chapter{Marco te\'orico}
\label{cap:marco-teorico}

Este cap\'itulo re\'une, de forma compacta, los conceptos y normas que
sustentan las decisiones t\'ecnicas y metodol\'ogicas descritas en el resto
del informe. No repite el detalle de aplicaci\'on de cada concepto a
BIOPET (eso corresponde a los cap\'itulos~\ref{cap:requisitos} a
\ref{cap:despliegue}); se limita a definir el marco conceptual usado y a
se\~nalar d\'onde se aplica.

\section{Ingenier\'ia de requisitos}

La especificaci\'on de requisitos de BIOPET sigue la estructura de
ISO/IEC/IEEE~29148:2018, est\'andar internacional que define el contenido
m\'inimo de un documento de Especificaci\'on de Requisitos de Software
(SRS): descripci\'on global del producto, requisitos funcionales y no
funcionales enunciados de forma verificable, y trazabilidad expl\'icita
hacia artefactos derivados (historias de usuario, casos de uso, pruebas).
La redacci\'on individual de cada requisito sigue adem\'as las
caracter\'isticas de calidad de la \emph{INCOSE Guide to Writing
Requirements}~\cite{incose2023requirements} --- en particular
\emph{Unambiguous} (un requisito admite una \'unica interpretaci\'on) y
\emph{Verifiable} (existe un m\'etodo objetivo para confirmar su
cumplimiento) --- aplicadas en el cap\'itulo~\ref{cap:requisitos}. Las
historias de usuario y los casos de uso son las dos t\'ecnicas
complementarias de elicitaci\'on/especificaci\'on usadas para descomponer
cada requisito en un flujo de interacci\'on concreto, seg\'un la pr\'actica
comparada por Pacheco, Garc\'ia \& Reyes~\cite{pacheco2018requirements}
(cap\'itulo~\ref{cap:trabajos-relacionados}).

\section{Arquitectura de software: el modelo C4}

El modelo C4 de Brown~\cite{c4-model} describe la arquitectura de un
sistema en niveles de abstracci\'on progresivos (Contexto, Contenedores,
Componentes, C\'odigo), cada uno dirigido a una audiencia distinta. BIOPET
documenta formalmente los niveles de Contenedores y Componentes
(cap\'itulo~\ref{cap:arquitectura}, secciones~\ref{sec:c4-nivel-2} y
\ref{sec:c4-nivel-3}), eligiendo ese nivel de detalle porque es el que
permite razonar sobre las decisiones de despliegue y de dise\~no interno
del backend sin caer en el detalle de clases individuales, que cambia con
mayor frecuencia que la arquitectura.

\section{Calidad de software: ISO/IEC~25010}

ISO/IEC~25010:2011~\cite{iso25010-2011} define un modelo de calidad de
producto de software organizado en ocho caracter\'isticas (Funcionalidad,
Eficiencia de desempe\~no, Compatibilidad, Usabilidad, Confiabilidad,
Seguridad, Mantenibilidad, Portabilidad), cada una con subcaracter\'isticas
espec\'ificas. Se usa en este informe como vocabulario com\'un para
organizar la evidencia emp\'irica recolectada (cap\'itulo~\ref{cap:pruebas-calidad}),
siguiendo el enfoque de aplicaci\'on pr\'actica descrito por Estdale \&
Georgiadou~\cite{estdale2018applying25010}: mapear cada caracter\'istica
contra evidencia real, declarando expl\'icitamente cu\'ando una
subcaracter\'istica no fue evaluada, en vez de omitirla silenciosamente
(secci\'on~\ref{sec:iso25010}).

\section{Arquitectura REST y JWT}

La API de BIOPET sigue el estilo arquitect\'onico REST (\emph{Representational
State Transfer}) sobre HTTP: recursos identificados por URL, verbos HTTP
con sem\'antica definida, y respuestas sin estado del lado del servidor
entre solicitudes. La autenticaci\'on usa JSON Web Token (JWT), formato
compacto y autocontenido definido en RFC~7519~\cite{rfc7519}: un token
firmado digitalmente que transporta un conjunto de \emph{claims}
(declaraciones) verificables sin necesidad de una consulta adicional al
servidor. BIOPET usa JWT firmado con HMAC-SHA256, transportado en cookies
en vez de en el cuerpo de la respuesta o en \texttt{localStorage}
(cap\'itulo~\ref{cap:seguridad}), decisi\'on de arquitectura registrada en
ADR-006. Los errores de la API siguen RFC~7807~\cite{rfc7807}
(\emph{Problem Details for HTTP APIs}), formato est\'andar para comunicar
fallos HTTP de forma estructurada y legible por m\'aquina.

\section{Seguridad de aplicaciones web: OWASP Top~10}

El OWASP Top~10~\cite{owasp-top10} es la clasificaci\'on de referencia de
la industria para los riesgos de seguridad m\'as cr\'iticos en aplicaciones
web (control de acceso roto, fallas criptogr\'aficas, inyecci\'on,
configuraci\'on de seguridad, fallas de identificaci\'on/autenticaci\'on,
fallas de registro/monitoreo, entre otros). BIOPET usa esta clasificaci\'on
como marco de auditor\'ia manual (cap\'itulo~\ref{cap:seguridad}),
complementada con an\'alisis din\'amico (OWASP ZAP) y est\'atico (SpotBugs
+ Find Security Bugs), siguiendo una l\'ogica de defensa en profundidad:
ning\'un m\'etodo \'unico de auditor\'ia se considera suficiente por s\'i
solo.

\section{Acceso a datos: JPA y procedimientos almacenados}

Spring Data JPA (\emph{Java Persistence API}) provee un mapeo
objeto-relacional (ORM) que traduce operaciones sobre entidades Java a
sentencias SQL generadas autom\'aticamente, adecuado para operaciones CRUD
simples sobre una \'unica tabla. Para operaciones que involucran l\'ogica
de negocio, validaci\'on cruzada entre tablas o agregaciones complejas,
BIOPET usa procedimientos y funciones almacenados de PostgreSQL, invocados
desde el backend mediante \texttt{@Procedure}/\texttt{@NamedStoredProcedureQuery}
de Spring Data JPA (secci\'on~\ref{sec:acceso-jpa-formal}). Esta
estrategia h\'ibrida sigue la misma l\'ogica documentada por Zmaranda et
al.~\cite{zmaranda2020crud}: dejar que cada mecanismo de acceso a datos
resuelva el tipo de operaci\'on para el que fue dise\~nado, en vez de
forzar toda la l\'ogica a trav\'es de un \'unico mecanismo. Esta decisi\'on
es consistente con evidencia emp\'irica comparativa adicional: SQL nativo
supera en tiempo de ejecuci\'on a las consultas generadas por ORM tanto en
C\#~\cite{nowicki2022comparative} como, v\'ia JDBC directo, en
Java~\cite{zuchnik2021comparative}, y la implementaci\'on de la capa de
persistencia tiene un impacto medible sobre el rendimiento general de la
aplicaci\'on~\cite{siebyla2020impact}. La comparaci\'on de rendimiento
entre pilas REST equivalentes (Spring~Boot y .NET~Core), relevante para la
decisi\'on de arquitectura registrada en ADR-002, sigue tambi\'en una
metodolog\'ia de \emph{benchmarking} comparable a la usada
aqu\'i~\cite{dhalla2021performance,godinho2024performance}.

\section{Patr\'on \emph{cache-aside}}

El patr\'on \emph{cache-aside} (tambi\'en llamado \emph{lazy loading})
consiste en que la aplicaci\'on consulta primero la cach\'e; si el dato no
est\'a presente (\emph{cache miss}), lo obtiene del almac\'en de datos
primario y lo escribe en la cach\'e para lecturas futuras; en cada
escritura sobre el dato primario, la entrada de cach\'e correspondiente se
invalida. BIOPET implementa este patr\'on de forma declarativa con las
anotaciones \texttt{@Cacheable} (lectura) y \texttt{@CacheEvict}
(invalidaci\'on) de Spring sobre Redis/Valkey
(secci\'on~\ref{sec:rendimiento-final}, listado~\ref{lst:cache-aside}),
respaldado por trabajos previos sobre la efectividad de esta estrategia en
acceso a datos relacionales~\cite{zulfa2020caching}. La configuraci\'on
del tiempo de vida (TTL) de cada entrada de cach\'e, externalizada por
variable de entorno en BIOPET (\texttt{CACHE\_TTL\_MS}), sigue la misma
recomendaci\'on de an\'alisis a gran escala de cl\'usteres de cach\'e en
memoria: la configuraci\'on del TTL es un factor clave del comportamiento
observado~\cite{yang2021large}. El despliegue contenerizado de BIOPET
(cap\'itulo~\ref{cap:despliegue}) tiene, como cualquier despliegue en
contenedores, un \emph{overhead} medible seg\'un su
configuraci\'on~\cite{ghatrehsamani2020art}.

\section{S\'intesis}

Los conceptos anteriores no se aplican de forma aislada: la ingenier\'ia
de requisitos (29148/INCOSE) alimenta la matriz de trazabilidad
(cap\'itulo~\ref{cap:trazabilidad-final}); el modelo C4 documenta la
arquitectura que implementa esos requisitos; ISO/IEC~25010 organiza la
evidencia emp\'irica que verifica si esa arquitectura cumple sus atributos
de calidad; y REST/JWT/OWASP/JPA/\emph{cache-aside} son las decisiones de
implementaci\'on concretas evaluadas por esa evidencia. Este marco
com\'un es lo que permite, en el cap\'itulo~\ref{cap:discusion}, discutir
los resultados de BIOPET en un mismo vocabulario en vez de como hallazgos
inconexos.
